% 148-14(148쪽) — 평행사변형 ABCD 와 대각선 AC · $\seg{AB}=2$, $\seg{BC}=4$, $B=60^\circ$.
% 스캔 화질이 낮아 TikZ 로 다시 그림. 원문과 같이 B 를 왼쪽 아래, C 를 오른쪽 아래, A·D 를 그 위에 두고 안쪽을 연하게 칠했다.
% 라벨은 원문에 있는 A·B·C·D·$2$·$4$·$60^\circ$ 만 둔다 — 넓이 $x$ 와 $\seg{AC}$ 의 길이 $y$(답)는 적지 않는다.
\begin{tikzpicture}[line width=0.5pt, scale=0.6, >=stealth]
  \coordinate (B) at (0,0);
  \coordinate (C) at (4,0);
  \coordinate (A) at (1,1.732);
  \coordinate (D) at (5,1.732);
  \fill[brown!14] (A) -- (B) -- (C) -- (D) -- cycle;
  \draw[line width=0.55pt] (A) -- (B) -- (C) -- (D) -- cycle;
  \draw[line width=0.55pt] (A) -- (C);
  \draw[line width=0.4pt] ([shift=(0:0.6)]B) arc (0:60:0.6);
  \node[font=\footnotesize] at ([shift=(30:1.02)]B) {$60^\circ$};
  \draw[densely dotted, line width=0.35pt, <->] ([shift=(150:0.2)]B) to[bend left=12]
    node[midway, left=-1.5pt, font=\footnotesize] {$2$} ([shift=(150:0.2)]A);
  \draw[densely dotted, line width=0.35pt, <->] ([shift=(-90:0.2)]B) to[bend right=12]
    node[midway, below=-2.5pt, font=\footnotesize] {$4$} ([shift=(-90:0.2)]C);
  \node[above left=-3.5pt and -3pt, font=\small] at (A) {$\pt{A}$};
  \node[left=1pt, font=\small] at (B) {$\pt{B}$};
  \node[below right=-3pt and -3pt, font=\small] at (C) {$\pt{C}$};
  \node[above right=-3.5pt and -3pt, font=\small] at (D) {$\pt{D}$};
\end{tikzpicture}
